
\section{Dyskusja}

Skutki niskiej jakości ocen organoleptycznych najlepiej opisano na przykładzie
win. W badaniach nad sędziami winiarskimi tylko około 10\% z nich
potrafiło powtórzyć własną ocenę tej samej, zaślepionej próbki wina
w obrębie jednej kategorii \cite{hodgson2008judge}. Co więcej, zdobycie
wysokiej nagrody na jednym konkursie nie pozwalało przewidzieć wyniku na
innym: oba wyniki były statystycznie niezależne \cite{hodgson2009concordance}. 
Ocenę organoleptyczną można skalibrować poprzez selekcję i~trening asesorów,
monitorowanie wskaźników jakości panelu (powtarzalności, odtwarzalności,
zgodności) oraz stosowanie wzorców odniesienia i~skal o~zewnętrznym
zakotwiczeniu \cite{sipos2021panel,bartoshuk2004valid}. Dla miodu metodykę
taką opracowano w~ramach International Honey Commission \cite{Piana2004},
a~we Włoszech od 1999~r. prowadzony jest krajowy rejestr ekspertów analizy
organoleptycznej miodu, do którego wpis wymaga ukończenia szkoleń i~zdania
egzaminu \cite{AlboMiele1999}. \textcolor{purple}{We Francji metodykę oceny
organoleptycznej miodów odmianowych rozwinął zespół M.~Gonneta (INRA,
obecnie INRAE), od którego prac wywodzi się współczesna analiza
organoleptyczna miodu \cite{Piana2004}. W~Niemczech zasady oceny
organoleptycznej miodu opracowało Niemieckie Towarzystwo Rolnicze (DLG):
jako jej cele wskazano m.in. charakterystykę miodów odmianowych łącznie
z~analizą pyłkową oraz wykrywanie skutków przegrzania i~długiego
przechowywania, a~schemat oceny przewiduje dla zapachu i~smaku m.in. wady
,,nietypowy dla miodu'', ,,obcy zapach'' i~,,fermentujący'' \cite{DLG2015}.
W~obu krajach ocena organoleptyczna uzupełnia badania laboratoryjne, a~nie
je zastępuje.} W~rutynowej ocenie
miodu wskaźniki te zwykle nie są jednak publikowane w~sprawozdaniach, a~laboratoria
nie dysponują wspólnym materiałem odniesienia. Porównywalność werdyktów
organoleptycznych między laboratoriami pozostaje więc założeniem, a~nie
wykazaną właściwością pomiaru.

Ten sam wzorzec widać między laboratoriami. Zgodność czterech laboratoriów
była słaba na skali binarnej i~nieznaczna na skali trzystopniowej
\cite{LandisKoch1977}, a~pełna zgodność dotyczyła wyłącznie próbek, do
których żadne laboratorium nie miało zastrzeżeń. Próbki P1 i~P2
zakwestionowało tylko jedno z~czterech laboratoriów, podczas gdy pozostałe
opisały je jako typowe, lekko typowe lub o~cechach niejednoznacznych.
Lab4, stosujące inną metodę, orzekło ,,cechy niejednoznaczne'' dla
wszystkich siedmiu próbek, także dla P6 i~P7, zakwestionowanych przez
pozostałe trzy laboratoria, i~opisało zapach charakterystyczny dla oznak
fermentacji w~trzech próbkach, w~tym w~P5, do której żadne z~pozostałych
laboratoriów nie miało zastrzeżeń. Werdykt
dla tej samej partii zależy więc nie tylko od losowej zmienności percepcji,
lecz także od przyjętej metody i~laboratorium.

Rozbieżności te nie dowodzą błędu żadnego konkretnego laboratorium ani wady
konkretnej partii. Dowodzą natomiast, że typowość zapachu i~smaku nie jest
w~praktyce komercyjnej kontroli jakości mierzona powtarzalnie między
laboratoriami -- zgodnie z~przewidywaniami literatury psychometrycznej
\cite{bartoshuk2003labeled,hodgson2009concordance,sipos2021panel}. Etykieta
,,nietypowy'' jest wynikiem porównania z~wewnętrznym, niejawnym wzorcem
oceniającego, a~ten różni się między osobami i~laboratoriami.

\subsection{Opis cechy a~jej waga decyzyjna}

Panel zarejestrował cechę obecną w~części opisów laboratoryjnych: słodycz
P6 i~P7 częściej niż innych próbek oceniano jako ,,za silną'', a~próbki te
uzyskały najniższą typowość. Zależność ta nie była jednak jednoznaczna.
Konsumenci przypisali taką samą intensywność słodyczy próbce~P5, do której
laboratoria stosujące kwalifikator typowości nie miały zastrzeżeń. Próbkę~P2,
opisaną przez jedno z~laboratoriów jako ,,bardzo słodka'', oceniali natomiast
jako zbyt słodką rzadko.

Wyniki te sugerują, że opis organoleptyczny może trafnie oddawać cechę
próbki, a~mimo to nie odróżniać miodu typowego od nietypowego. Można to
tłumaczyć tym, że o~typowości decyduje nie sama intensywność pojedynczej
cechy, lecz jej zgodność z~profilem danej odmiany. Wymaga to jednak
dalszych badań. Dlatego warto odróżnić \emph{trafny opis} od
\emph{rozstrzygającego kryterium}: cecha może być przydatna do
charakterystyki produktu, a~jednocześnie nieprzydatna jako samodzielna
przesłanka decyzji administracyjnej \cite{Gajda2026}.

\subsection{Konsekwencje dla kontroli jakości}

Z~punktu widzenia kontroli istotne jest nie to, czy ocena smaku zawiera
jakąkolwiek informację, lecz czy pozwala powtarzalnie odtworzyć werdykt.
Pojedyncza ocena panelisty odtwarzała werdykt większości laboratoriów na
poziomie $\mathrm{AUC} = 0{,}57$ (95\% CI: 0{,}54--0{,}60), a~reguła oparta
na typowości smaku myliłaby się w~blisko co drugiej ocenie. Zgodność
konsumenta z~laboratorium ($\kappa = 0{,}10$) nie różniła się od zgodności
dwóch konsumentów ($\kappa = 0{,}08$). Same laboratoria zgadzały się w~pełni
co do trzech z~siedmiu próbek na skali binarnej i~jednej na skali
trzystopniowej. Kryterium, którego wynik dla próbek granicznych zależy od
wyboru oceniającego lub laboratorium, nie spełnia warunku powtarzalności,
a~użyte samodzielnie -- warunku proporcjonalności.

Różnicę tę widać w~samych sprawozdaniach. Lab2 i~Lab4 opisywały
regułę podejmowania decyzji z~uwzględnieniem niepewności pomiaru
\cite[pkt~7.8.6]{ISO17025} w~odniesieniu do wartości granicznych parametrów
mierzalnych, a~Lab4 wskazało wprost parametry fizykochemiczne
i~analizę pyłkową. Dla oceny organoleptycznej żadne sprawozdanie nie
podawało niepewności ani reguły decyzyjnej. W~przypadku metod
instrumentalnych precyzję wyznacza się w~badaniach międzylaboratoryjnych,
a~na jej podstawie ustala różnice krytyczne stosowane przy ocenie zgodności
z~kryteriami czystości miodu \cite{Ulberth2024}. Dla werdyktu ,,smak
nietypowy'' analogiczne wielkości nie są znane.

Nie oznacza to, że ocena organoleptyczna jest bezużyteczna. Jest cennym
narzędziem badania preferencji, opisu i~doskonalenia produktu, a~także
sygnałem do dalszych badań, w~roli przewidzianej dla testu organoleptycznego
w~skoordynowanym planie kontroli Komisji Europejskiej
\cite{KE2015PlanKontroli}, gdy wynik jest skrajny i~zgodny między
niezależnymi oceniającymi -- jak dla P6 i~P7. Podobną zależność opisano dla
wyszkolonych ekspertów: w~próbach International Honey Commission wyniki
oceny zgodności z~profilem miodu odmianowego były ważne dla próbek wyraźnie
dobrych lub wadliwych, natomiast dla próbek o~cechach pośrednich rozrzut
ocen był największy, a~binarna klasyfikacja ,,tak/nie'' nie dawała ważnego
wyniku w~23\% przypadków nawet po połączeniu trzech paneli
\cite{Piana2004}. \textcolor{purple}{Wyniki wskazują, że przy próbkach
granicznych werdykt organoleptyczny nie powinien być potwierdzany kolejną
oceną organoleptyczną prowadzoną bez wzorca. Drugie laboratorium,
posługujące się własnymi deskryptorami, dostarcza kolejnego opisu, a~nie
rozstrzygnięcia: jedno może określić smak jako ,,bardzo słodki'', inne
zapach jako ,,landrynkowy'', i~bez zdefiniowanego wzorca deskryptorów
typowości żaden z~tych opisów nie przesądza o~nietypowości. Werdykt ,,smak
nietypowy'' powinien zatem zostać zweryfikowany wynikami miarodajnymi:
oznaczeniami fizykochemicznymi i~mikroskopowymi, a~przy podejrzeniu
zafałszowania metodami analitycznymi o~znanej precyzji \cite{Ulberth2024},
zanim stanie się podstawą decyzji o~wycofaniu produktu. Rola oceny
organoleptycznej w~kontroli jakości handlowej powinna ograniczać się do
wykrywania cech nietypowych określonych w~przepisach, takich jak obcy smak
lub zapach oraz oznaki fermentacji
\cite{Dyrektywa2001110,RozporzadzenieMRiRW2009}, a~nie obejmować orzekania
o~typowości, dla której wzorzec nie istnieje.}

\textcolor{purple}{Takie ograniczenie odpowiada konstrukcji prawa unijnego.
Dyrektywa 2001/110/WE formułuje wymagania organoleptyczne w~postaci
zakazów: miód, z~wyjątkiem miodu piekarskiego, nie może mieć obcego smaku
ani zapachu i~nie może wykazywać oznak fermentacji \cite{Dyrektywa2001110}.
Ocena wyglądu, zapachu, smaku, konsystencji, krystalizacji i~widocznych
oznak fermentacji pozwala więc wskazać możliwą niezgodność i~uzasadnić
skierowanie próbki do dalszych badań laboratoryjnych. W~tej roli,
wstępnej i~uzupełniającej, a~nie jako samodzielną metodę analityczną,
wykorzystano test organoleptyczny w~skoordynowanym planie kontroli
Komisji Europejskiej \cite{KE2015PlanKontroli}.}

\subsection{Ograniczenia badania}
\label{sub:ograniczenia}

Najważniejszym ograniczeniem jest mała liczba partii. Choć panel wystawił
966~ocen, porównania na poziomie próbek (korelacja panel--laboratoria,
zgodność werdyktów) opierają się na siedmiu jednostkach. Po drugie, panel tworzyli
konsumenci po krótkim szkoleniu, a~nie wyselekcjonowani asesorzy; oddaje
to perspektywę rynkową, lecz zwiększa rozrzut ocen. Pytanie panelu
o~typowość smaku nie jest też tożsame z~procedurami laboratoryjnymi, które
oceniają łącznie zapach i~smak według własnych wzorców. Po trzecie, ocenę
laboratoryjną przeprowadzono jako zlecenia komercyjne, a~nie jako formalny
test pierścieniowy: próbki nie były kodowane, laboratoria nie dysponowały
wspólnym materiałem odniesienia, a~kodowanie ich opisów na skali T/L/N,
w~tym przypisanie werdyktowi ,,cechy niejednoznaczne'' kodu L, jest
uproszczeniem dokonanym przez autorów. Po czwarte, panel (maj--lipiec 2026)
i~ocenę laboratoryjną (sierpień 2026) dzieliło od jednego do trzech miesięcy,
w~czasie których cechy zależne od przechowywania (krystalizacja, fermentacja) mogły
się zmienić. Wreszcie, badano tylko siedem partii, dlatego wyników nie należy
uogólniać na wszystkie miody.

\section{Wnioski}

Ocena typowości smaku i~zapachu miodu prowadzona tak jak w~tym
badaniu (przez ślepy panel konsumencki lub przez pojedyncze
laboratorium, bez wspólnego wzorca odniesienia i~bez raportowanej miary
zgodności oceniających) jest zbyt zróżnicowana, by mogła być
rozstrzygającym parametrem oceny jakości handlowej próbki.

Przemawiają za tym trzy niezależne ustalenia. Po pierwsze, zdolność
dyskryminacyjna oceny jest nieakceptowalnie niska: pojedyncza ocena
panelisty odtwarzała werdykt laboratoriów z~$\mathrm{AUC} = 0{,}57$, czyli
na poziomie niewiele lepszym od wyboru binarnie losowej wartości ($\mathrm{AUC} = 0{,}5$)
i~wyraźnie poniżej wartości 0{,}7, uznawanej za minimum akceptowalnej
dyskryminacji \cite{Hosmer2013}. W~praktyce reguła ,,smak nietypowy --
próbka odrzucona'' zakwestionowałaby 39{,}7\% ocen próbek, do których
większość laboratoriów nie miała zastrzeżeń, i~przepuściła 46{,}4\% ocen
próbek przez nią zakwestionowanych, myliłaby się więc w~blisko co drugiej
decyzji. Po drugie, poza jedną wyraźnie typową próbką panel nie różnicował
próbek, a~jako całość nie zakwestionował żadnej z~nich: typowość próbek
zakwestionowanych przez większość laboratoriów (0{,}438 i~0{,}489) była
nieodróżnialna od 0{,}5. Po trzecie, oceniający nie byli zgodni ani między
sobą, ani z~laboratoriami: zgodność czterech laboratoriów była słaba na
skali binarnej ($\kappa = 0{,}30$) i~nieznaczna na skali trzystopniowej
($\kappa = 0{,}13$), a~pojedynczy konsument zgadzał się z~laboratorium
w~podobnym stopniu jak z~innym konsumentem ($\kappa = 0{,}10$
i~$\kappa = 0{,}08$).

Decyzja o~uznaniu partii za zgodną lub o~jej wycofaniu, oparta wyłącznie na
takim wyniku, zależałaby w~większym stopniu od tego, kto i~gdzie ocenia, niż
od właściwości samej próbki. Badanie subiektywne nie powinno więc stanowić
podstawy binarnej oceny „dobry/zły''. Ocena organoleptyczna zachowuje wartość
jako narzędzie opisowe i~sygnał do dalszych badań, zwłaszcza gdy wynik jest
skrajny i~zgodny między niezależnymi oceniającymi. Nie powinna jednak
samodzielnie rozstrzygać o~zgodności miodu z~wymaganiami jakości handlowej.
\textcolor{purple}{Ocena organoleptyczna powinna ograniczać się do
wykrywania cech nietypowych określonych w~przepisach, takich jak obcy smak
lub zapach oraz oznaki fermentacji, a~nie orzekać o~typowości, dla której
nie ma wzorca. Werdykt ,,smak nietypowy'' powinien być każdorazowo
weryfikowany wynikami miarodajnymi, a~nie kolejną oceną organoleptyczną
prowadzoną bez wzorca deskryptorów.}
\textcolor{purple}{W~konsekwencji decyzja o~uznaniu miodu za niezgodny
z~wymaganiami wyłącznie na podstawie nieokreślonego wrażenia
organoleptycznego, np.\ ,,zbyt słodki'', ,,niedostatecznie aromatyczny''
lub ,,niedostatecznie typowy'', byłaby podatna na zakwestionowanie, jeżeli
takiego kryterium nie określa prawo Unii, przepisy krajowe, specyfikacja
produktu ani uznana lub zwalidowana metoda badań
\cite[art.~34]{Rozporzadzenie2017625}.}
